\appendix

\section{核心模块代码}

由于代码量较大，本附录只列出最能反映流水线设计核心思想的几个模块：mips.v（顶层连接与flushE合成）、datapath.v（数据通路）、hazard.v（冒险处理）。其余基础模块（alu、mux、触发器等）与实验三类似，限于篇幅不再赘述。

\subsection{mips.v（顶层连接）}

\begin{lstlisting}[language=Verilog,basicstyle=\ttfamily\footnotesize,breaklines=true]
module mips(
    input  wire        clk, rst,
    input  wire [31:0] instr,
    output wire [31:0] pc,
    input  wire [31:0] readdata,
    output wire        memwrite,
    output wire [31:0] aluout,
    output wire [31:0] writedata
);
    wire [5:0]  opD, functD;
    wire        jumpD, branchE;
    wire        regdstE, alusrcE;
    wire [2:0]  alucontrolE;
    wire        memwriteM, memtoregM, regwriteM;
    wire        memtoregW, regwriteW, memtoregE;
    wire [4:0]  rsD, rtD, rsE, rtE;
    wire [4:0]  writeregE, writeregM, writeregW;
    wire [1:0]  forwardAE, forwardBE;
    wire        stallF, stallD;
    wire        flushE_hazard, flushD, flushE_total;

    // 关键合成：beq/j成立时，ID/EX 必须清零
    assign flushE_total = flushE_hazard | flushD;

    controller u_ctrl(.clk(clk), .rst(rst), .opD(opD), .functD(functD),
        .flushE(flushE_total), .jumpD(jumpD), .branchE(branchE),
        .regdstE(regdstE), .alusrcE(alusrcE), .alucontrolE(alucontrolE),
        .memtoregE(memtoregE), .memwriteM(memwriteM), .memtoregM(memtoregM),
        .regwriteM(regwriteM), .memtoregW(memtoregW), .regwriteW(regwriteW));

    datapath u_dp(.clk(clk), .rst(rst), .instrF(instr), .readdataM(readdata),
        .jumpD(jumpD), .branchE(branchE), .regdstE(regdstE), .alusrcE(alusrcE),
        .alucontrolE(alucontrolE), .memtoregW(memtoregW), .regwriteW(regwriteW),
        .forwardAE(forwardAE), .forwardBE(forwardBE),
        .stallF(stallF), .stallD(stallD), .flushE(flushE_total),
        .opD(opD), .functD(functD), .rsD(rsD), .rtD(rtD),
        .rsE(rsE), .rtE(rtE), .writeregE(writeregE),
        .writeregM(writeregM), .writeregW(writeregW),
        .pcF(pc), .aluoutM(aluout), .writedataM(writedata),
        .flushD_out(flushD));

    hazard u_hz(.rsD(rsD), .rtD(rtD), .rsE(rsE), .rtE(rtE),
        .writeregE(writeregE), .memtoregE(memtoregE),
        .writeregM(writeregM), .regwriteM(regwriteM),
        .writeregW(writeregW), .regwriteW(regwriteW),
        .forwardAE(forwardAE), .forwardBE(forwardBE),
        .stallF(stallF), .stallD(stallD), .flushE(flushE_hazard));

    assign memwrite = memwriteM;
endmodule
\end{lstlisting}

\subsection{hazard.v（冒险处理）}

\begin{lstlisting}[language=Verilog,basicstyle=\ttfamily\footnotesize,breaklines=true]
module hazard(
    input  wire [4:0] rsD, rtD, rsE, rtE,
    input  wire [4:0] writeregE, writeregM, writeregW,
    input  wire       memtoregE, regwriteM, regwriteW,
    output wire [1:0] forwardAE, forwardBE,
    output wire       stallF, stallD, flushE
);
    // EX 前推 A 路：MEM 优先于 WB
    assign forwardAE =
        ((rsE != 5'b0) && (rsE == writeregM) && regwriteM) ? 2'b10 :
        ((rsE != 5'b0) && (rsE == writeregW) && regwriteW) ? 2'b01 : 2'b00;
    // EX 前推 B 路
    assign forwardBE =
        ((rtE != 5'b0) && (rtE == writeregM) && regwriteM) ? 2'b10 :
        ((rtE != 5'b0) && (rtE == writeregW) && regwriteW) ? 2'b01 : 2'b00;
    // Load-use 暂停
    wire lwstall;
    assign lwstall = memtoregE && ((rsD==rtE)||(rtD==rtE)) && (rtE!=5'b0);
    assign stallF = lwstall;
    assign stallD = lwstall;
    assign flushE = lwstall;
endmodule
\end{lstlisting}

\subsection{datapath.v 关键片段}

\begin{lstlisting}[language=Verilog,basicstyle=\ttfamily\footnotesize,breaklines=true]
// IF 阶段：PC 与 PC+4
flopenr #(32) u_pcreg(.clk(clk), .rst(rst), .en(~stallF),
    .d(pcnextF), .q(pcF));
adder u_pcplus4(.a(pcF), .b(32'd4), .y(pcplus4F));
mux2 #(32) u_branch_mux(.d0(pcplus4F), .d1(pcbranchE),
    .sel(pcsrcE), .y(pc_seq_or_branch));
mux2 #(32) u_jump_mux(.d0(pc_seq_or_branch), .d1(pcjumpD),
    .sel(jumpD), .y(pcnextF));

// IF/ID 寄存器：flushD 由分支或跳转触发
assign flushD = pcsrcE | jumpD;
assign flushD_out = flushD;
flopenrc #(32) u_ifid_instr(.clk(clk), .rst(rst),
    .en(~stallD), .clear(flushD), .d(instrF), .q(instrD));

// ID/EX 寄存器：clear = flushE（顶层合成的flushE_total）
floprc #(32) u_idex_rd1(.clk(clk), .rst(rst), .clear(flushE),
    .d(rd1D), .q(rd1E));
// ... 其他ID/EX寄存器类似

// EX 阶段：三路前推与ALU
mux3 #(32) u_fwdA(.d0(rd1E), .d1(resultW), .d2(aluoutM),
    .sel(forwardAE), .y(srcaE));
mux3 #(32) u_fwdB(.d0(rd2E), .d1(resultW), .d2(aluoutM),
    .sel(forwardBE), .y(srcb2E));
mux2 #(32) u_alusrc(.d0(srcb2E), .d1(signimmE),
    .sel(alusrcE), .y(srcbE));
alu u_alu(.a(srcaE), .b(srcbE), .alucontrol(alucontrolE),
    .result(aluoutE), .zero(zeroE));

// beq 判断在 EX
assign pcsrcE = branchE & zeroE;

// sw 写入数据使用前推后的 srcb2E
flopr #(32) u_exmem_wd(.clk(clk), .rst(rst),
    .d(srcb2E), .q(writedataM));
\end{lstlisting}

\section{累加测试程序（汇编与机器码）}

\begin{lstlisting}[basicstyle=\ttfamily\footnotesize]
# 地址  机器码      汇编                  注释
0x00:  20010000   addi $1, $0, 0       # sum = 0
0x04:  20020001   addi $2, $0, 1       # i = 1
0x08:  20030064   addi $3, $0, 100     # limit = 100
0x0C:  00220820   add  $1, $1, $2      # sum += i  <- loop:
0x10:  20420001   addi $2, $2, 1       # i++
0x14:  0062202a   slt  $4, $3, $2      # $4 = (100 < i)
0x18:  1080fffc   beq  $4, $0, loop    # if (i <= 100) goto loop
0x1C:  ac010000   sw   $1, 0($0)       # mem[0] = sum (5050)
0x20:  08000008   j    8                # end: 死循环
\end{lstlisting}

\section{Testbench 代码}

\begin{lstlisting}[language=Verilog,basicstyle=\ttfamily\footnotesize,breaklines=true]
`timescale 1ns / 1ps
module testbench();
    reg clk, rst;
    wire [31:0] writedata, dataadr, display_data;
    wire memwrite;

    top dut(.clk(clk), .rst(rst), .writedata(writedata),
        .dataadr(dataadr), .memwrite(memwrite),
        .display_data(display_data));

    initial begin rst = 1; #22; rst = 0; end
    initial begin clk = 0; forever #10 clk = ~clk; end

    always @(negedge clk) begin
        if (memwrite) begin
            if (dataadr === 32'd0 && writedata === 32'd5050) begin
                $display("================================================");
                $display(" Simulation succeeded ");
                $display(" sum(1+2+...+100) = %0d (0x%h)",
                         writedata, writedata);
                $display("================================================");
                $stop;
            end else begin
                $display(" Simulation FAILED ");
                $display(" dataadr=%0d writedata=%0d", dataadr, writedata);
                $stop;
            end
        end
    end
endmodule
\end{lstlisting}
